\scanpage{74}
Reformulation des données lorsque $c\notin C$ en ttes fibres
i.e. $\rho$ inv ($n$ impair, $\rho'$ inv, ou $n$ pair, $\rho'$ inv. et $2$ inv.).

A) $\varphi_G:G\longrightarrow\operatorname{Aut}(X)$
ou encore $(\tau_i^X)_{0\leq i\leq n}$

est remplacé par
\[
\varphi_G:G\longrightarrow\operatorname{Aut}(V)
\quad\text{ou encore }(\tau_i^V)_{0\leq i\leq n}
\]

Conditions

$1^\circ)$ Les $\tau_i^V$ pseudo-réflexions
\[
(\text{\unsure{déf}} \quad L_i\subset\unsure{V_i|V},\ H_i\subset\unsure{V_i|V})
\]

\margin{\deleted{\unsure{dans les fibres}}}

$2^\circ)$
\[
V=\bigoplus_{0\leq i\leq n}L_i,\quad
\check V\simeq\bigoplus\underbrace{H_i^0}_{\unsure{L'_0|L'_i}}
\quad\text{i.e. }\bigcap_{0\leq i\leq n}H_i=0
\quad\text{\unsure{ch. q. fibre}}
\]

$3^\circ)$ $L_i\subset H_{i+1}$ ch. q. fibre
\hfill$0\leq i\leq n-1$

\deleted{$4^\circ)$ Posons $D_0=\bigcap_{1\leq i\leq n}H_i$,
\unsure{i.e. le}}

B) La donnée de $s_0\in\unsure{\Gamma(X)|X}$
\deleted{$\mathbin{-}$}
donnant lieu avec $\varphi_G$ à une orb. géom non dég,
est \deleted{\unsure{un pt}} exprimée par
\[
\uline{\widetilde s_0\text{ base de}}\quad
D_0=\bigcap_{1\leq i\leq n}H_i
\]
qui épingle la situation.
\deleted{On choisit}
[Il revient au $=$ de se donner une base
\deleted{C) Les formes quadratiques}
$a'_0$ de $L'_0$ (car $D_0\otimes L'_0\xrightarrow{\sim}\mathcal O_S$, par
\[
\boxed{\langle\widetilde s_0,a'_0\rangle=1}
\]
ce qui fixe \unsure{base} $a_0$ de $L_0$ par
$\tau_0=\operatorname{id}_V+a'_0\otimes a_0$,
et de proche en proche on détermine les $a'_i,a_i$ par
\[
\langle a_i,a'_{i+1}\rangle=1,\qquad
\tau_i=\operatorname{id}+a'_i\otimes a_i
\quad]
\]
\mspage{73}

